\section{Integrating the Web Project}
\label{sec:web}

\subsection{NuGet Package Changes}

Two additions are required in the web project's \texttt{.csproj}: a reference to \texttt{Aspire.Npgsql.EntityFrameworkCore.PostgreSQL} and a project reference to the ServiceDefaults library:

\begin{lstlisting}[language=XML]
<!-- src/Game.Web/Game.Web.csproj (additions only) -->
<ItemGroup>
  <PackageReference Include="Aspire.Npgsql.EntityFrameworkCore.PostgreSQL"
                    Version="9.0.0" />
</ItemGroup>

<ItemGroup>
  <ProjectReference Include="..\Game.ServiceDefaults\
Game.ServiceDefaults.csproj" />
</ItemGroup>
\end{lstlisting}

\texttt{Aspire.Npgsql.EntityFrameworkCore.PostgreSQL} is a thin integration package that wraps \texttt{Npgsql.EntityFrameworkCore.PostgreSQL} and adds Aspire-aware configuration reading. Specifically, its \texttt{AddNpgsqlDbContext\textless TContext\textgreater} extension method reads the connection string from the .NET configuration system under the key \texttt{ConnectionStrings\_\_\{name\}}, which is exactly the key that Aspire injects via \texttt{WithReference}. The developer does not write a connection string anywhere; Aspire provides it at runtime.

\subsection{Program.cs Changes}

Two additions are required in the web project's \texttt{Program.cs}. Both must appear before \texttt{builder.Build()}:

\begin{lstlisting}[language={[Sharp]C}]
// src/Game.Web/Program.cs (additions only; existing lines unchanged)
var builder = WebApplication.CreateBuilder(args);

// --- Aspire additions ---
builder.AddServiceDefaults();                             // (1)
builder.AddNpgsqlDbContext<AppDbContext>("gamedb");        // (2)
// --- end Aspire additions ---

// ... existing DI registrations, middleware, etc. ...

var app = builder.Build();

// --- Aspire addition ---
app.MapDefaultEndpoints();                                // (3)
// --- end Aspire addition ---

// ... existing middleware pipeline, app.Run() ...
\end{lstlisting}

Three annotations:

\textbf{(1) \texttt{AddServiceDefaults()}.} Registers OpenTelemetry, HTTP resilience, and service discovery. Must be called before any \texttt{AddHttpClient} calls so that the resilience and service-discovery handlers are applied globally.

\textbf{(2) \texttt{AddNpgsqlDbContext\textless AppDbContext\textgreater("gamedb")}.} Registers \texttt{AppDbContext} in the DI container, reads the connection string from \texttt{ConnectionStrings\_\_gamedb}, and enables Aspire health checks and telemetry for the database connection. The string \texttt{"gamedb"} must match the argument passed to \texttt{AddDatabase} in the AppHost. No \texttt{appsettings.json} entry is needed.

\textbf{(3) \texttt{MapDefaultEndpoints()}.} Registers \texttt{/health} and \texttt{/alive} endpoints. This call must appear after \texttt{builder.Build()} and before \texttt{app.Run()}.

\subsection{Avoiding Conflicts with Existing DbContext Registrations}

If the web project already registers \texttt{AppDbContext} via \texttt{AddDbContext} (for example, with a connection string read from \texttt{appsettings.json}), that registration must be removed before adding the Aspire version. Registering the same context twice produces a runtime \texttt{InvalidOperationException}. The Aspire-aware registration supersedes the manual one and does not require any compensating configuration entry in \texttt{appsettings.json}.

\subsection{Database Initialisation at Startup}

For the Game POC, schema creation uses \texttt{EnsureCreatedAsync}

\begin{lstlisting}[language={[Sharp]C}]
// Inside app.Services.CreateScope() at startup
await dbContext.Database.EnsureCreatedAsync(); // POC / development
// OR
await dbContext.Database.MigrateAsync();       // production / staged migrations
\end{lstlisting}

Both calls are safe to make under \texttt{WaitFor} ordering because Aspire guarantees the Postgres container is accepting connections before the web project process starts.
